\section{工程实现与代码结构}

本章介绍项目的工程实现情况, 包括项目代码结构、核心模块说明、训练评估流程和遇到的工程问题.

\subsection{项目目录结构}

本项目采用模块化的代码组织结构, 核心目录如下:

\begin{lstlisting}[language=bash,caption={项目目录结构}]
project_root/
  src/
    data/          # 数据加载、预处理、CLIP 缓存
    models/        # 模型定义: CLIP, A2 EEG encoder, fusion, caption model
    train/         # 训练脚本: baseline, fusion, alignment, semantic fusion
    eval/          # 评估: generation, retrieval, sanity, metrics, gate analysis
    losses/        # 损失函数: contrastive, EEG augmentations, similarity
    utils/         # 工具: config, seed, checkpoint, logger
  scripts/         # Shell/Python 启动脚本
  configs/         # YAML 配置文件
  outputs/         # 训练输出、checkpoint、metrics、报告
  tests/           # 单元测试
\end{lstlisting}

\subsection{核心代码模块}

表~\ref{tab:code-modules} 列出了项目的主要代码模块及其功能.

\begin{table}[H]
  \centering
  \caption{核心代码模块说明}
  \label{tab:code-modules}
  \begin{tabularx}{\linewidth}{l X}
    \toprule
    模块路径 & 功能说明 \\
    \midrule
    \texttt{src/data/dataset.py} & EEGVisionCaptionDataset: 支持 JSONL manifest 的 EEG-Image-Caption 数据集加载器, 支持 Thought2Text、EEG-ImageNet、THINGS-EEG2 等数据格式 \\
    \texttt{src/data/clip\_cache.py} & CLIP 视觉特征预缓存: 全量图像 pooled embedding 和 visual tokens 的预计算与存储 \\
    \texttt{src/models/eeg\_encoder.py} & 13 种 EEG encoder variant（含 A2 TemporalSpectralSpatialTransformer）、工厂函数 build\_eeg\_encoder \\
    \texttt{src/models/fusion.py} & GatedFusion: 门控残差融合模块, 支持 image\_only fallback \\
    \texttt{src/models/caption\_model.py} & SoftPromptCaptionModel: 从 512-d embedding 到软提示 token 的投影 + frozen LM. 含内置 byte-level tiny LM 回退 \\
    \texttt{src/models/alignment\_model.py} & EEGCLIPAlignmentModel: EEG→CLIP 对齐训练模型 \\
    \texttt{src/train/train\_fusion.py} & 融合模型训练: image + EEG→gated fusion→caption CE loss \\
    \texttt{src/train/train\_align.py} & EEG-to-CLIP 对齐训练: InfoNCE + MSE + optional class CE \\
    \texttt{src/train/train\_semantic\_fusion.py} & A2 语义融合训练 \\
    \texttt{src/eval/generate.py} & 生成式评估: 多种 epilepsy 模式下跑 full-test \\
    \texttt{src/eval/sanity\_check.py} & 四种 EEG 模式 sanity check \\
    \texttt{src/eval/metrics.py} & 指标统计, JSONL 输出 \\
    \texttt{src/eval/semantic\_fusion\_classifier\_eval.py} & A2 semantic fusion 分类评估 \\
    \texttt{src/utils/config.py} & YAML 配置加载（含简单 YAML 回退解析器, 无需 pyyaml）\\
    \texttt{src/utils/seed.py} & 全局随机种子设置 \\
    \bottomrule
  \end{tabularx}
\end{table}

\subsection{训练与评估流程}

\textbf{训练流程}:

\begin{enumerate}
  \item 数据准备: \texttt{python scripts/make\_dummy\_data.py}（调试）/ 运行 \texttt{build\_thought2text\_manifest.sh}（正式数据）;
  \item CLIP 预缓存: \texttt{python scripts/precompute\_vision.py} 预计算 visual features 和 text prototypes;
  \item A2 对齐训练（可选）: \texttt{python -m src.train.train\_align --config configs/base.yaml};
  \item A2 语义融合训练: \texttt{python -m src.train.train\_semantic\_fusion};
  \item VTF 训练（需要 enhanced tokens）: 通过 token-level 训练脚本;
  \item 生成式 EVLM 训练: \texttt{python -m src.train.train\_clip\_adapter} 或相应 LoRA 训练脚本, 使用 Qwen2-VL;
  \item 评估: \texttt{python -m src.eval.generate --mode real\_eeg / vision\_only / shuffled\_eeg / random\_eeg}.
\end{enumerate}

\textbf{输出文件}: 每次训练在 \texttt{outputs/<run\_name>/} 下生成 checkpoint（\texttt{best.pt, checkpoint\_last.pt}）、training log、metrics JSON 和 report Markdown. 生成式评估在 \texttt{eval/} 子目录下产出 JSONL 文件, 每条记录包含 \texttt{image\_id, mode, reference, prediction}.

\subsection{工程问题与解决}

\textbf{LoRA 依赖}: Qwen2-VL 需要 \texttt{transformers >= 4.45} 和 \texttt{peft} 库. 由于 Qwen2-VL 的视觉 adapter 与标准 CLIP ViT 的 token 接口不完全一致, 需要在 token fusion 后插入一个轻量级 adapter 进行格式对齐.

\textbf{显存控制}: Qwen2-VL-2B 在 bfloat16 下约占用 4.5GB 显存, 加上 CLIP ViT、A2 encoder、LoRA 参数和 batch 数据, 总显存约 6.5GB 到 13.7GB 不等（取决于 route 和 LoRA 配置）. 所有模型均在 48GB 单 GPU 上完成, 无需 4-bit 量化.

\textbf{Full-test evaluation}: 测试集共计 1997 trial × 6 退化 × 4 模式 ≈ 48,000 次前向推理. 生成式模型的自回归解码是主要瓶颈, 每个 caption 生成约需 0.2--0.5 秒. 全量测试约需 2--4 小时.

\textbf{Caption target 清洗}: BLIP 生成的伪 caption 中约 15\% 存在格式不规范（如包含特殊符号、过短或过长）. 对 BLIP captions 施加最低 3 词过滤, 并将过长的 caption（$>40$ token）截断. Invalid caption 过滤在评估阶段使用正则表达式自动检测 URL-like 输出、wnid 代码和重复 token.
